% Part: incompleteness
% Chapter: representability-in-q
% Section: representable-comp

\documentclass[../../../include/open-logic-section]{subfiles}

\begin{document}

\olfileid{inc}{req}{rpc}
\olsection{$\Th{Q}$માં નિરૂપ્ય વિધેયો સંગણનીય છે}

$\Th{Q}$માં નિરૂપ્ય દરેક વિધેય
સંગણનીય છે તે સાબિત કરીશું.
પહેલાં~$\Th{Q}$માં નિરૂપ્ય વિધેયો
વિશે એક લેમ્મા સ્થાપવો પડશે.

\begin{lem}\ollabel{lem:rep-q}
  $f(x_0,
  \dots, x_k)$ વિધેય~$\Th{Q}$માં
  નિરૂપ્ય હોય, તો એવું
  !!a{formula}~$!A_f(x_0, \dots, x_k, y)$
  છે કે
  \sourcecorrection{OLINC-022}{સ્થિર અંગ્રેજી
  મૂળના આ વાક્યમાં સૂત્રનું નામ
  $!A$ લખાયું છે, પરંતુ અગાઉની
  વ્યાખ્યા અને આ લેમ્માના તરત
  પછીના વિધાનમાં તે $!A_f$ છે.
  અહીં એ જ નામ સતત રાખ્યું છે.}
  \[
  \Th{Q} \Proves !A_f(\num{n_0}, \dots, \num{n_k}, \num{m})
  \quad\text{ત્યારે અને માત્ર ત્યારે}\quad m = f(n_0, \dots, n_k).
  \]
\end{lem}

\begin{proof}
  ``જો'' દિશા
  \olref[int]{defn:representable-fn}\olref[int]{defn:rep:a}
  પરથી મળે છે. ``માત્ર જો'' દિશા
  આ પ્રમાણે મળે: માનો કે
  $\Th{Q} \Proves !A_f(\num{n_0},
  \dots, \num{n_k}, \num{m})$,
  પરંતુ $m \neq f(n_0, \dots, n_k)$.
  $l = f(n_0, \dots, n_k)$ લો.
  \olref[int]{defn:representable-fn}\olref[int]{defn:rep:a}
  પ્રમાણે $\Th{Q}
  \Proves !A_f(\num{n_0}, \dots,
  \num{n_k}, \num{l})$.
  \olref[int]{defn:representable-fn}\olref[int]{defn:rep:b}
  પ્રમાણે
  $\lforall[y][(!A_f(\num{n_0}, \dots,
  \num{n_k}, y) \lif \num{l} =
  y)]$. તર્કના નિયમો અને
  $\Th{Q} \Proves
  !A_f(\num{n_0}, \dots,
  \num{n_k}, \num{m})$ની ધારણા વડે
  $\Th{Q} \Proves
  \eq[\num{l}][\num{m}]$ મળે છે.
  બીજી તરફ,
  \olref[bre]{lem:q-proves-neq} પ્રમાણે
  $\Th{Q} \Proves
  \eq/[\num{l}][\num{m}]$.
  એટલે $\Th{Q}$ અસંગત થાય.
  પરંતુ એવું શક્ય નથી, કારણ કે
  પ્રમાણભૂત નિદર્શ $\Th{Q}$ને સંતોષે
  છે (જુઓ
  \olref[int][def]{def:standard-model}),
  $\Sat{N}{\Th{Q}}$; અને યથાર્થતા
  પ્રમેય પ્રમાણે સંતોષી શકાય એવો
  સિદ્ધાંત હંમેશાં સુસંગત હોય છે
  (\tagrefs{prfAX/{fol:axd:sou:cor:consistency-soundness},
  prfSC/{fol:seq:sou:cor:consistency-soundness},
  prfND/{fol:ntd:sou:cor:consistency-soundness},
  prfTab/{fol:tab:sou:cor:consistency-soundness}}).
\end{proof}

\begin{lem}
$\Th{Q}$માં નિરૂપ્ય દરેક વિધેય
સંગણનીય છે.
\end{lem}

\begin{proof}
આ દાવો સાચો કેમ છે તેનો પહેલાં
સહજ વિચાર જોઈએ. $f$ ગણવા માટે
આ પ્રમાણે કરીએ. અંકગણિતની
ભાષામાં શક્ય બધી
!!{derivation}s~$\delta$ની યાદી
બનાવીએ. આ યાંત્રિક રીતે શક્ય છે.
દરેક માટે તપાસીએ કે તે
$!A_f(\num{n_0}, \dots,
\num{n_k}, \num{m})$ સ્વરૂપના
!!a{formula}ની !!a{derivation}
છે કે નહીં ($\Th{Q}$માં~$f$નું
નિરૂપણ કરતું આ !!{formula}
\olref{lem:rep-q}માંથી મળે છે).
એવું હોય, તો \olref{lem:rep-q}
પ્રમાણે $m = f(n_0, \dots, n_k)$
અને~$f$નું મૂલ્ય મળી ગયું.
શોધ પૂરી થાય છે, કારણ કે
$\Th{Q} \Proves
!A_f(\num{n_0}, \dots,
\num{n_k}, \num{f(n_0, \dots,
n_k)})$; તેથી આખરે યોગ્ય
$\delta$ મળી જ રહેશે.

આ દલીલ હજી સંપૂર્ણ ચોક્કસ નથી:
આપણી પ્રક્રિયા માત્ર સંખ્યાઓને
બદલે !!{derivation}s અને
!!{formula}s પર કામ કરે છે;
વળી ``શક્ય બધી !!{derivation}sની
યાદી'' યાંત્રિક રીતે બનાવી
શકાય તે પણ હજી સમજાવ્યું નથી.
પરંતુ પદો, !!{formula}s અને
!!{derivation}sને ગડેલ-સંખ્યાઓ
આપી શકાય છે તે આપણે જોયું છે.
સંગણનાનું ચોક્કસ નિદર્શ—સામાન્ય
પુનરાવર્તી વિધેયો—પણ રજૂ કર્યું છે.
અને $\Prf[\Th{Q}](d,y)$ સંબંધ,
જે ત્યારે અને માત્ર ત્યારે સત્ય છે
જ્યારે $d$ એ~$\Th{Q}$નાં
સ્વયંસિદ્ધ વાક્યોમાંથી ગડેલ-સંખ્યા~$y$
ધરાવતા !!{formula}ની
!!a{derivation}ની ગડેલ-સંખ્યા
હોય, (આદિમ) પુનરાવર્તી છે તે પણ
જોયું છે. આગળ જરૂરી અન્ય આદિમ
પુનરાવર્તી વિધેયો $\fn{num}$
(\olref[art][trm]{prop:num-primrec})
અને $\fn{Subst}$
(\olref[art][sub]{prop:subst-primrec}) છે.
આમાંથી લઘુત્તમીકરણ વડે~$f$
વ્યાખ્યાયિત કરી શકાય; તેથી~$f$
પુનરાવર્તી છે.

પહેલાં આ વ્યાખ્યા આપો:
\begin{multline*}
  A(n_0, \dots, n_k, m) = \\
  \fn{Subst}(\fn{Subst}(\dots\fn{Subst}(\Gn{!A_f}, \fn{num}(n_0), \Gn{x_0}),\\ \dots),
  \fn{num}(n_k),  \Gn{x_k}), \fn{num}(m), \Gn{y})
\end{multline*}
આ જટિલ દેખાય છે, પરંતુ તે માત્ર
$A(n_0, \dots, n_k,
m) = \Gn{!A_f(\num{n_0}, \dots,
\num{n_k}, \num{m})}$ વિધેય છે.

હવે $R(n_0, \dots, n_k, s)$
સંબંધ લો. $(s)_0$ એ
$\Th{Q}$માંથી
$!A_f(\num{n_0}, \dots,
\num{n_k}, \num{(s)_1})$ની
!!a{derivation}ની ગડેલ-સંખ્યા
હોય ત્યારે આ સંબંધ સત્ય છે:
\[
R(n_0, \dots, n_k, s) \quad\text{ત્યારે અને માત્ર ત્યારે}\quad \Prf[\Th{Q}]((s)_0, A(n_0,
\dots, n_k, (s)_1))
\]
$R(n_0, \dots, n_k, s)$ સત્ય થાય
એવું~$s$ શોધીએ, તો સંખ્યાઓની
જોડ~$(s)_0$ અને~$(s)_1$ મળે છે,
જેમાં $(s)_0$ એ
$A_f(\num{n_0}, \dots,
\num{n_k}, \num{(s)_1})$ની
!!a{derivation}ની ગડેલ-સંખ્યા છે.
\sourcecorrection{OLINC-023}{સ્થિર અંગ્રેજી મૂળના
આ ગદ્યમાં છેલ્લી દલીલ
$(s)_1$ લખાઈ છે, જ્યારે તેના
પહેલાંનું વિધાન અને સંબંધ બંનેમાં
સૂત્રની અંદર અંકનિરૂપણ
$\num{(s)_1}$ જરૂરી છે.
અહીં ગદ્યને એ જ સૂત્ર સાથે
મેળમાં કર્યું છે.}
આમ~$s$ શોધવું એ અનૌપચારિક
સાબિતીમાં~$d$ અને~$m$ની જોડ
શોધવા જેવું છે. એવું~$s$
``શોધતું'' સંગણનીય વિધેય નિયમિત
લઘુત્તમીકરણ વડે વ્યાખ્યાયિત થાય.
યાદ રાખો કે $R$ નિયમિત છે:
દરેક $n_0$, \dots, $n_k$
માટે $\Th{Q} \Proves
!A_f(\num{n_0}, \dots,
\num{n_k}, \num{f(n_0, \dots,
n_k)})$નું કોઈ
!!a{derivation}~$\delta$ છે;
તેથી $s = \tuple{\Gn{\delta},
f(n_0, \dots, n_k)}$ માટે
$R(n_0, \dots, n_k, s)$ સત્ય છે.
એટલે $f$ને આ પ્રમાણે લખી શકીએ:
\[
f(n_0,\dots,n_{k}) = (\umin{s}{R(n_0, \dots, n_k, s)})_1.
\]
\end{proof}

\end{document}
